\documentclass[12pt,a4paper]{article}

%=========================================================
%                   ENCODAGE ET POLICES
%=========================================================
\usepackage[utf8]{inputenc}
\usepackage[T1]{fontenc}
\usepackage[french]{babel}
\usepackage{lmodern}

%=========================================================
%                     MATHÉMATIQUES
%=========================================================
\usepackage{amsmath,amssymb,amsfonts,mathrsfs}
\usepackage{tkz-tab}
\everymath{\displaystyle}

%=========================================================
%                     MISE EN PAGE
%=========================================================
\usepackage[left=1.5cm,right=1.5cm,top=1.5cm,bottom=1.8cm]{geometry}
\usepackage{multicol}
\usepackage{float}
\usepackage{array,multirow}
\usepackage{enumitem}

%=========================================================
%                GRAPHISMES ET COULEURS
%=========================================================
\usepackage{tikz}
\usetikzlibrary{arrows.meta,calc}
\usepackage{pgfplots}
\pgfplotsset{compat=1.15}

\usepackage[most]{tcolorbox}
\tcbuselibrary{skins,hooks}

%=========================================================
%                  LIENS ET ARRIÈRE-PLAN
%=========================================================
\usepackage[colorlinks=true,linkcolor=blue,urlcolor=blue,citecolor=blue]{hyperref}
\usepackage{eso-pic}

%=========================================================
%                   COULEURS ET STYLES
%=========================================================
\definecolor{myred}{RGB}{180, 40, 40}
\definecolor{myblue}{RGB}{20, 50, 120}
\definecolor{mygreen}{RGB}{30, 120, 40}

%=========================================================
%                       FILIGRANE
%=========================================================
\AddToShipoutPicture{
    \AtTextCenter{
        \makebox[0pt]{
            \rotatebox{45}{
                \textcolor[gray]{0.92}{
                    \fontsize{5cm}{5cm}\selectfont PGB
                }
            }
        }
    }
}

\begin{document}

% En-tête
\begin{center}
    \Large\textbf{\underline{\textcolor{myred}{Suites Numériques}}}\\[0.2cm]
    \normalsize\textbf{Prof : M. BA} \hfill \textbf{Classe : Terminale L}\\[0.1cm]
    \textbf{Durée : 6 heures} \hfill \textbf{Année scolaire : 2025 -- 2026}
\end{center}
\hrule
\vspace{0.3cm}

%=========================================================
%    I. Généralités
%=========================================================
\section*{\color{myred}I. Généralités}

\subsection*{\color{myred}1. Activité introductive}

\begin{tcolorbox}[colback=blue!5!white,colframe=myblue,title=\textbf{Activité : Découverte de la notion de suite}]
On considère la liste de nombres réels suivants : $-10 \,;\, 5 \,;\, 3 \,;\, 4 \quad \text{et} \quad -8$.

\begin{enumerate}
    \item Ranger ces $5$ nombres dans l'ordre \textbf{décroissant}.
    \item En attribuant le **rang $0$** au plus grand nombre, donner le rang de chacun des autres nombres de la liste ordonnée.
\end{enumerate}
\end{tcolorbox}

\vspace{0.2cm}

\paragraph{\color{mygreen}Correction de l'activité :}
\begin{enumerate}
    \item L'ordre décroissant des $5$ nombres est : $5 \,;\, 4 \,;\, 3 \,;\, -8 \,;\, -10$.
    \item Attribution des rangs :
    \begin{itemize}[label=$\bullet$]
        \item Le rang de $5$ est : \textbf{0}
        \item Le rang de $4$ est : \textbf{1}
        \item Le rang de $3$ est : \textbf{2}
        \item Le rang de $-8$ est : \textbf{3}
        \item Le rang de $-10$ est : \textbf{4}
    \end{itemize}
\end{enumerate}

\vspace{0.3cm}


\begin{tcolorbox}[colback=yellow!5!white,colframe=myred,title=\textbf{Exploitation pédagogique}]
\begin{itemize}[label=$\Rightarrow$]
    \item La liste ordonnée $5 \,;\, 4 \,;\, 3 \,;\, -8 \,;\, -10$ est un exemple de **suite de nombres réels**.
    \item Si l'on nomme cette suite $u$, chaque nombre de la liste est un **terme** de la suite $u$, repéré de manière unique par son **rang** (ou indice).
    \item Pour noter un terme, on écrit le nom de la suite et on place son rang en indice :
    \[
    u_0 = 5 \quad;\quad u_1 = 4 \quad;\quad u_2 = 3 \quad;\quad u_3 = -8 \quad;\quad u_4 = -10
    \]
    \item De façon générale, le terme situé au rang $n$ est noté $\mathbf{u_n}$ (lu \textit{« u indice n »} ou \textit{« u de n »}).
\end{itemize}
\end{tcolorbox}

\subsection*{\color{myred}2. Définition}

\begin{tcolorbox}[colback=red!5!white,colframe=myred,title=\textbf{Définition}]
On appelle \textbf{suite numérique} toute fonction $u$ définie d'une partie $E$ de $\mathbb{N}$ (souvent $\mathbb{N}$ lui-même) vers $\mathbb{R}$.

\[
\begin{array}{cccc}
u : & \mathbb{N} & \to & \mathbb{R} \\
    & n & \mapsto & u_n
\end{array}
\]
Le réel $u_n$ est appelé \textbf{terme général} ou \textbf{terme d'indice $n$}. La suite est notée $(u_n)_{n \in \mathbb{N}}$ ou simplement $(u_n)$.
\end{tcolorbox}
\end{document}